\section{Empirical trends}\label{sec:discussion}
We end our discussion of experimental results by briefly reporting on several trends that we observed across multiple datasets.

\subsection{Underspecification}
Prior work has shown that there is often insufficient information at training time to distinguish models that would generalize well under distribution shift;
many models that perform similarly in-distribution (ID) can vary substantially out-of-distribution (OOD) \citep{mccoy2019berts,zhou2020curse,damour2020underspecification}.
In \Wilds, we attempt to alleviate this issue by providing multiple training domains in each dataset as well as an OOD validation set for model selection.
Perhaps as a result, we do not observe significantly higher variance in OOD performance than ID performance in \reftab{erm_drops}, with the exception of
\Amazon and \CivilComments, where the OOD performance is measured on a smaller subpopulation and is therefore naturally more variable.
Excluding those datasets, the average standard deviation from \reftab{erm_drops} is 2.6\% for OOD performance and 2.0\% for ID performance, which is comparable.
These results raise the question of when underspecification, as reported in prior work, could be more of an issue.

\subsection{Model selection with in-distribution versus out-of-distribution validation sets}
All of the baseline results reported in this paper use an OOD validation set for model selection, as discussed in \refsec{erm_drops_model_selection}.
To facilitate research into comparisons of ID versus OOD performance, most \Wilds datasets also provide an ID validation and/or test set.
For example, in \iWildCam, the ID validation set comprises photos from the same set of camera traps used for the training set.
These ID sets are not used for model selection nor official evaluation.

\citet{gulrajani2020search} showed that on the DomainBed domain generalization datasets, selecting models with an ID validation set leads to higher OOD performance than using an OOD validation set.
This contrasts with our approach of using OOD validation sets, which we find to generally provide a good estimate of OOD test performance.
Specifically, in \refapp{experiments_id_vs_ood}, we show that for our baseline models, model selection using an OOD validation set results in comparable or higher OOD performance than model selection using an ID validation set.
This difference could stem from many factors: for example, \Wilds datasets tend to have many more domains, whereas DomainBed datasets tend to have fewer domains that can be quite different from each other (e.g., cartoons vs. photos);
and there are some differences in the exact procedures for comparing performance using ID versus OOD validation sets.
Further study of the effects of these different model selection procedures and choices of validation sets would be a useful direction for future work.

\subsection{The compounding effects of multiple distribution shifts}
Several \Wilds datasets consider hybrid settings, where the goal is to simultaneously generalize to unseen domains as well as to certain subpopulations.
We observe that combining these types of shifts can exacerbate performance drops.
For example, in \PovertyMap and \FMoW, the shift to unseen domains exacerbates the gap in subpopulation performance (and vice versa).
Notably, in \FMoW, the difference in subpopulation performance (across regions) is not even manifested until also considering another shift (across time).
While we do not always observe the compounding effect of distribution shifts---e.g., in \Amazon, subpopulation performance is similar whether we consider shifts to unseen users or not---these observations underscore the importance of evaluating models on the combination of distribution shifts that would occur in practice, instead of considering each shift in isolation.
